% =============================================================================
%  Routing Results -- Brick routing results on Dataset A.
%  Include-ready LaTeX section: no documentclass, no preamble, no bibliography.
% =============================================================================

\section{Routing Results}
\label{sec:routing-results}

This section reports routing-system results on Dataset~A. The present revision
discusses only the Brick router. Additional routing baselines can be added to
this file in later revisions, but they are intentionally omitted here to keep
the interpretation focused on Brick's calibrated preference profiles.

\subsection{Brick Experimental Protocol}
\label{sec:routers-brick-protocol}

Brick is evaluated on all \(5{,}504\) Dataset~A rows using the locked production
configuration described in the companion algorithm section. The reported
numbers use stratified 5-fold out-of-fold skill-vector calibration:
for each fold, the router's per-model capability table is calibrated on the
other four folds and evaluated on the held-out fold. The routing hyperparameters
are not re-selected inside this full-dataset evaluation; they are fixed from
prior development.

This protocol uses the entire dataset while avoiding per-example leakage in the
skill-vector estimates. It should be read as a calibrated in-domain evaluation,
not as a zero-shot result and not as validation on an independent external
benchmark. No target model weights are updated.

\subsection{Brick Preference Profiles}
\label{sec:routers-brick-profiles}

The Brick router exposes a continuous routing preference \(r\in[-1,1]\). We
evaluate five operating points:
\[
  r\in\{-1.0,-0.5,0.0,+0.5,+1.0\}.
\]
The endpoints correspond to named user profiles:
\[
  \mathrm{min}\equiv -1.0,\qquad
  \mathrm{neutral}\equiv 0.0,\qquad
  \mathrm{max}\equiv +1.0.
\]
The \texttt{min} profile prioritizes cost reduction, the \texttt{neutral}
profile targets the calibrated quality-price compromise, and the \texttt{max}
profile prioritizes answer quality by strongly relaxing cost pressure.

\subsection{Full-Dataset Results}
\label{sec:routers-brick-results}

Table~\ref{tab:brick-routing-oof-results} reports Brick's full-dataset
out-of-fold routing curve. The primary quality metric is selected-answer
accuracy: a routed query is counted correct if the model selected by the router
solves the query. Route-exact accuracy is reported separately; it measures
agreement with the cheapest-correct label and therefore penalizes intentional
over-routing to stronger models.

\begin{table}[ht]
\centering
\caption{Brick routing results on all \(5{,}504\) Dataset~A rows using locked
production hyperparameters and stratified 5-fold out-of-fold skill calibration.
The distribution column reports selected-model percentages as
qwen/deepseek/kimi.}
\label{tab:brick-routing-oof-results}
\small
\begin{tabular}{l r r r r l}
\toprule
\textbf{Profile} & \textbf{\(r\)} & \textbf{Selected answer} &
\textbf{95\% CI} & \textbf{Avg. cost} & \textbf{Distribution} \\
\midrule
min     & -1.0 & 63.17\% & \(\pm 1.27\) pp & 0.100 & 100/0/0 \\
low     & -0.5 & 63.17\% & \(\pm 1.27\) pp & 0.100 & 100/0/0 \\
neutral & \phantom{-}0.0 & 68.46\% & \(\pm 1.23\) pp & 0.192 & 71/25/3 \\
high    & +0.5 & 74.18\% & \(\pm 1.16\) pp & 0.359 & 25/59/16 \\
max     & +1.0 & \textbf{76.98\%} & \(\pm 1.11\) pp & 0.537 & 0/31/69 \\
\bottomrule
\end{tabular}
\end{table}

\begin{table}[ht]
\centering
\caption{Route-exact accuracy for the same Brick profiles. Route-exact accuracy
matches the cheapest-correct label, so it is aligned with economy routing but
not with the max-performance objective.}
\label{tab:brick-routing-route-exact}
\small
\begin{tabular}{l r r}
\toprule
\textbf{Profile} & \textbf{\(r\)} & \textbf{Route-exact accuracy} \\
\midrule
min     & -1.0 & 63.17\% \\
low     & -0.5 & 63.17\% \\
neutral & \phantom{-}0.0 & 60.10\% \\
high    & +0.5 & 37.39\% \\
max     & +1.0 & 17.84\% \\
\bottomrule
\end{tabular}
\end{table}

Table~\ref{tab:brick-routing-baseline-context} gives the single-model and oracle
context for interpreting Table~\ref{tab:brick-routing-oof-results}. These rows
are not routing algorithms; they are reference points for the attainable
response-accuracy range within the three-model pool.

\begin{table}[ht]
\centering
\caption{Full-dataset context baselines for Dataset~A. The oracle row counts a
query as solved if any model in the pool solves it.}
\label{tab:brick-routing-baseline-context}
\small
\begin{tabular}{l r r}
\toprule
\textbf{Reference} & \textbf{Selected answer} & \textbf{Avg. cost} \\
\midrule
always \texttt{qwen3.5-9b}        & 63.17\% & 0.100 \\
always \texttt{deepseek-v4-flash} & 73.69\% & 0.400 \\
always \texttt{kimi2.6}           & 75.02\% & 0.600 \\
three-model oracle                & 83.25\% & n/a \\
\bottomrule
\end{tabular}
\end{table}

\subsection{Discussion}
\label{sec:routers-brick-discussion}

The \texttt{min} profile collapses to the cheapest-model baseline: it selects
\texttt{qwen3.5-9b} for every query, obtaining \(63.17\%\) selected-answer
accuracy at normalized cost \(0.100\). This is the intended behaviour for the
lowest-cost setting: the router suppresses over-capacity and cost-insensitive
escalation.

The \texttt{neutral} profile increases selected-answer accuracy to \(68.46\%\),
a \(+5.29\) percentage-point gain over \texttt{min}, while increasing average
cost from \(0.100\) to \(0.192\). Operationally, this is the quality-price
compromise profile: it still routes \(71\%\) of queries to the cheapest model,
but uses \texttt{deepseek-v4-flash} on roughly one quarter of the corpus and
\texttt{kimi2.6} on a small tail of high-requirement queries.

The \texttt{high} and \texttt{max} profiles show the intended aggressive scaling
of the preference knob. At \texttt{high}, Brick reaches \(74.18\%\), close to
the best single-model baselines, while still selecting \texttt{qwen3.5-9b} on
about one quarter of the dataset. At \texttt{max}, Brick reaches
\(\mathbf{76.98\%}\), exceeding the best single-model baseline,
\texttt{kimi2.6} at \(75.02\%\), by \(+1.96\) percentage points. The remaining
gap to the \(83.25\%\) three-model oracle is the residual routing opportunity
not captured by the current feature representation and decision rule.

The route-exact curve should be interpreted differently from selected-answer
accuracy. Route-exact accuracy falls from \(63.17\%\) at \texttt{min} to
\(17.84\%\) at \texttt{max}, not because the max profile fails to answer, but
because it intentionally chooses stronger models even when a cheaper model would
also be correct. For performance-first routing, selected-answer accuracy is the
appropriate metric; for economy-first routing, route-exact accuracy is a useful
diagnostic of cheapest-correct matching.
